\appendix
\section*{Appendix}

\section{Proof of Lemma \ref{lemma:bijection}}

{\em Proof.}
We write $R|_{X \times Y}$ for $\{ (x,y) \in X \times Y \ |\ xRy \}$.

Let $R'$ be $R|_{\dom(R) \times \dom(S)}$.

Claim 1. $\dom(R')=\dom(R)$.
We can show it as follows.
Fix $x \in \dom(R)$.
Since $R,S$ is domain-stable, we have $\dom(RS) = \dom(R)$.
Hence $x \in \dom(RS)$.
Hence there are $y$ and $z$ such that $x R y S z$.
Hence $y \in \dom(S)$.
Hence $x \in \dom(R')$.

Claim 2. $\codom(R')=\dom(S)$.
We can show it as follows.
By definition of $R'$, we have $\codom(R') \subseteq \dom(S)$.
Since $R'$ is injective, we have $|\dom(R')| \le |\codom(R')|$.
From Claim 1, we have $\dom(R')=\dom(R)$.
Hence $|\dom(R)| \le |\codom(R')| \le |\dom(S)|$.
Since $|\dom(R)| = |\dom(S)|$, we have $\codom(R')=\dom(S)$.

$R'^{-1}$ is a function $\dom(S) \to \dom(R)$.
We can show it as follows.
Fix $x \in \dom(S)$ to show $x R'^{-1} y$ for some $y$.
By Claim 2. $x \in \codom(R')$. Hence there is $y$ such that $y R' x$.
Hence $x R'^{-1} y$.
Assume $x R'^{-1} y$ and $x R'^{-1} z$ to show $y=z$.
Then $y R' x$ and $z R' x$.
Since $R'$ injective, we have $y=z$.

The function $R'^{-1}$ is surjective.
We can show it as follows.
Fix $x \in \dom(R)$.
By Claim 1, we have $x \in \dom(R')$.
Hence there is $y$ such that $x R' y$.
Hence $y R'^{-1} x$.

Since $|\dom(S)|=|\dom(R)|$, the function $R'^{-1}$ is bijective.
Hence $R'$ is a bijection.
Hence $R:\dom(R) \to \dom(S)$ is a bijection.
$\Box$

\section{Proof of Lemma \ref{lemma:shorten}}

{\em Proof.}

Assume $|\pi| > C_0$
to show that there are $\pi'$ and $\sigma'$ such that
they are obtained from $\pi$ and $\sigma$ by replacing some subsequences
and $W(\pi,\sigma) = W(\pi',\sigma')$ and $|\pi|>|\pi'|$.

Let $\pi$ be $(R_1,\ldots,R_n)$.
Let $\sigma$ be $(\Vec v_1,\ldots,\Vec v_n)$.
Let $k=|\dom(R_n)|$.
Let $W(\pi,\sigma)=(W_1,\ldots,W_K)$.
Since $\pi$ is domain-stable, we have $W_i=(\infty,0)$ for all $i \in [k+1,K]$.

Take $p_0 \in [1,n]$ such that
$p_0$ is the smallest among $p_0$'s such that
$|\dom(R_{p_0})|>0$
Let $k_0$ be $|\dom(R_{p_0})|$.

Take $p_1 \in [p_0+1,n]$ such that 
$p_1$ is the smallest among $p_1$'s such that
$|\dom(R_{p_1})|>k_0$.

Case 1. $p_1-p_0 \ge |V|{K \choose k_0}+2$.

By Lemma \ref{lemma:maxlen}, we have
$\Val(R_{p_0},\Vec v_{p_0}) \ne \Val(R_{p_1-1},\Vec v_{p_1-1})$
and
$\Val(R_{p_0+1},\Vec v_{p_0+1}) \ne \Val(R_{p_1-1},\Vec v_{p_1-1})$.

Define $\pi'$ from $\pi$ by 
replacing $R_{p_0},R_{p_0+1}$ by a single composition $R_{p_0}R_{p_0+1}$.

Case 2. $p_1-p_0 \le |V|{K \choose k_0}+1$.

Since $C_0 = \Sigma_{k=1}^K|V|{K \choose k}+1$,
by the condition for length
there are some 
$k' \in [1,k]$
and
$p_2 \in [p_1,n]$
and
$p_3 \in [p_2+1,n]$
such that
for all $R \in \pi$ we have
$|\dom(R)|=k'$ iff $R=R_i$ for some $i \in [p_2,p_3-1]$
and
$p_3-p_2 > |V|{K \choose k'}$.

By Lemma \ref{lemma:maxlen},
$\Val(R_{p_2},\Vec v_{p_2}) \ne \Val(R_{p_3-1},\Vec v_{p_3-1})$.

Define $\pi'$ from $\pi$ by 
replacing $R_{p_2-1},R_{p_2}$ by a single composition $R_{p_2-1}R_{p_2}$.

In both cases, the following holds.

Define $\sigma'$ as the sequence obtained from $\sigma$
corresponding to $\pi'$.
Then $|\pi'|<|\pi|$ and $\pi'$ is domain-stable.
Moreover $\sigma'$ is along $\pi'$, by Lemma \ref{lemma:A}.

Then by comparing the $i$-th element for each $i \in [1,K]$,
we have
$W(\pi',\sigma')=W(\pi,\sigma)$.
$\Box$

\section{Proof of Proposition \ref{prop:1}}

{\em Proof.}
Assume $(c,x) >_\Pi (d,y)$.
There is a path $\pi=\pi_1,\ldots,\pi_k$ such that
$(c,x) >_\pi (d,y)$.
For $i \in [1,k]$, let $c_i$ and $t_i$ be
the companion number and the stage number in the bottom sequent
of the path $\pi_i$.
Let $c_{k+1}$ and $t_{k+1}$ be $d$ and $y$.
Then $(c_i,t_i)$ and $(c_{i+1},t_{i+1})$ are
the companion numbers and the stage numbers in the bottom sequent
and the top sequent of the path $\pi_i$.

By Proposition \ref{prop:0.5},
$\rho'(c_i,t_i) >_\Lex \rho'(c_{i+1},t_{i+1})$ for $i \in [1,k]$.
Hence
$\rho(c,x) >_\Lex \rho'(d,y)$
$\Box$

\section{Proof of Proposition \ref{prop:2}}


{\em Proof.}
Assume $\forall cx(\forall dy(\rho'(d,y) <_U \rho'(c,x) \imp Gdy) \imp Gcx)$
to show $\forall cx.Gcx$ (labeled with (c)).

Let B n be $\forall cx(\rho'(c,x) \le n \imp Gcx)$.

We will use the following induction principle for $>_\Lex$:
\[
\forall n((\forall m<_U n.Bm) \imp Bn) \imp \forall n.Bn.
\]

We will show
\[
\forall n((\forall m<_U n.Bm) \imp Bn)
\]
(labeled with (b)).
Fix $n$ and 
assume $\forall m<_U n.Bm$ to show $Bn$.
Fix $c,x$ and assume $\rho'(c,x) \le_U n$ to show $Gcx$ (labeled with (a)).

Let $\Pi_1$ be the subtree of root $c$. Let $c_1,\ldots,c_k$ be the leaves
of $\Pi_1$.
For $i \in [1,k]$,
let $\pi_i$ be the path from $c$ to $c_i$.

Let $v_i$ be the stage-number of the sequent $c_i$
when the stage number of $c$ is $x$ in the subproof $\Pi_1$.
By the ranking function, we have
$\rho(c,((x)_0,\ldots,(x)_{l-1})) >_\Lex
\rho(d,((v_i)_0,\ldots,(v_i)_{l_i-1}))$
where $l =|x|$ and $l_i=|v_i|$.
Hence
$\rho'(c,x) >_\Lex \rho'(d,v_i)$.

From the assumption $\rho'(c,x) \le_U n$, we have
$\rho'(d,y) <_U n$.
In the assumption $\forall m<_U n.Bm$,
by taking $m$ to be $\rho'(c_i,v_i)$,
we have $B(\rho'(c_i,v_i))$.
In the definition of $B$, by taking $c$ and $x$ to be $c_i$ and $v_i$,
we have $Gc_iv_i$. Hence $G_{c_i}v_i$.Hence
\[
(\Gamma_{2c_i})^\circ_{(v_i)_0\ldots (v_i)_{l_i-1}} \prove \Delta_{2c_i}^\bullet
\]
where $\Gamma_{2c_i}$ is the bottom sequent of the path $\pi_{c_i}$ and
where $\Gamma_{1c_i}$ is the top sequent of the path $\pi_{c_i}$.

Hence every open assumption of $\Pi_1$ is provable.
Hence the conclusion
$(\Gamma_{2c})^\circ_{(x)_0\ldots (x)_{l-1}} \prove \Delta_{2c}^\bullet$
is provable.
Hence $Gcx$ holds, which shows (a).
Hence  we haves shown (b).
By induction principle for $>_\Lex$ we have $\forall n.Bn$.

We will show (c).
Fix $c,x$.
In $\forall n.Bn$, by taking $n$ to be $\rho'(c,x)$, we have $Gcx$.
Hence we have shown (c).
$\Box$

\section{Inference Rules for LJ}

The inference rules for LJ is given in Figure \ref{fig:1a}.

\begin{figure*}[t]
{\prooflineskip
\[
\infer[(\Cont)]{\Gamma,A \prove \Delta}{\Gamma,A,A \prove \Delta}
\qquad
\infer[(\Exch)]{\Gamma,A,B,\Gamma' \prove \Delta}{\Gamma,B,A,\Gamma' \prove \Delta}
\\
\infer[(\Axiom)]{\Gamma,A \prove A,\Delta}{}
\qquad
\infer[(\Wk)]{\Gamma,A \prove \Delta}{\Gamma \prove \Delta}
%\\
\qquad
\infer[(\Cut)]{\Gamma \prove \Delta}{
\Gamma \prove F,\Delta
&
\Gamma,F \prove \Delta
}
%\qquad
\\
\infer[(\Subst)]{\Gamma\theta \prove \Delta\theta}{\Gamma \prove \Delta}
\qquad
\infer[(\neg L)]{\Gamma,\neg F \prove \Delta}{\Gamma \prove F,\Delta}
%\\
\qquad
\infer[(\neg R)]{\Gamma \prove \neg F,\Delta}{\Gamma,F \prove \Delta}
\qquad
\infer[(\lor L)]{\Gamma,F \lor G \prove \Delta}{
\Gamma,F \prove \Delta
&
\Gamma,G \prove \Delta
}
%\qquad
\\
\infer[(\lor R)]{\Gamma \prove F \lor G,\Delta}{
	\Gamma \prove F,G,\Delta
}
\qquad
\infer[(\land L)]{\Gamma,F \land G \prove \Delta}{\Gamma,F,G \prove \Delta}
\qquad
\infer[(\land R)]{\Gamma \prove F \land G,\Delta}{
\Gamma \prove F,\Delta
&
\Gamma \prove G,\Delta
}
\\
\infer[(\imp L)]{\Gamma, F \imp G \prove \Delta}{
	\Gamma \prove F,\Delta
	&
	\Gamma,G \prove \Delta
}
\qquad
\infer[(\imp R)]{\Gamma \prove F \imp G,\Delta}{\Gamma,F \prove G,\Delta}
\qquad
\infer[(\forall L)]{\Gamma,\forall xF \prove \Delta}{\Gamma,F[x:=t] \prove \Delta}
\\
%\]
%\[
\infer[(\forall R)]{\Gamma \prove \forall xF,\Delta}{\Gamma \prove F,\Delta}
(x \notin \FV(\Gamma))
\qquad
\infer[(\exists L)]{\Gamma,\exists xF \prove \Delta}{\Gamma,F \prove \Delta}
(x \notin \FV(\Gamma,\Delta))
%\\
\qquad
\infer[(\exists R)]{\Gamma \prove \exists xF,\Delta}{\Gamma \prove F[x:=t],\Delta}
%\qquad
\\
\infer[(=L)]{\Gamma[x:=u],t=u \prove \Delta[x:=u]}{\Gamma[x:=t] \prove \Delta[x:=t]}
\qquad
\infer[(=R)]{\Gamma \prove t=t,\Delta}{}
\]
}
\vskip -3ex
\caption{Inference Rules}\label{fig:1a}
\vskip -1ex
\end{figure*}

\section{Example of Cyclic Proofs}

\begin{Eg}\label{eg:1}\rm
An example of a cyclic proof is given in Figure \ref{fig:5},
\begin{figure*}[t]
\[
\infer[(\lor L)]{Ex \lor Ox \prove \N x}{
	\infer[(\Case\ E)]{(a)\ Ex \prove \N x}{
		\infer[(=L)]{x=0 \prove \N x}{
		\infer[(Wk)]{x=0 \prove \N 0}{
		\infer[(\N\ R1)]{\prove \N 0}{}
		}}
		&
		\infer[(=L)]{x=sx',Ox' \prove \N x}{
		\infer[(\N\ R2)]{x=sx',Ox' \prove \N sx'}{
		\infer[(Wk)]{x=sx',Ox' \prove \N x'}{
		\infer[(\Subst)]{Ox' \prove \N x'}{
		(b)\ Ox \prove \N x
		}}}}
	}
	&
	\infer[(\Case\ O)]{(b)\ Ox \prove \N x}{
	\infer[(=L)]{x=sx',Ex' \prove \N x}{
	\infer[(\N\ R2)]{x=sx',Ex' \prove \N sx'}{
	\infer[(Wk)]{x=sx',Ex' \prove \N x'}{
	\infer[(\Subst)]{Ex' \prove \N x'}{
	(a)\ Ex \prove \N x
	}}}}}
}
\]
\caption{Cyclic Proof}\label{fig:5}
\end{figure*}
where the mark (a) and (b) denote the bud-companion relation, and
the production rules are
\[
\infer{E0}{}
\qquad
\infer{Esx}{Ox}
\qquad
\infer{Osx}{Ex}
\]
\end{Eg}
In Figure \ref{fig:5} the companion (a) uses $E x$, but
the bud (a) uses $E x$ where $x$ is $x'$ and $x' < x$,
so 
their actual meanings are different
even though they are of the same form.
The global trace condition generalizes this fact
and guarantees the soundness of a cyclic proof system.

\section{Addition of Heyting Arithmetic}

This subsection adds arithmetic to the logical systems.

We add arithmetic to both $\LJID$ and $\CLJIDomega$.
\begin{Def}\rm
$\CLJIDHA$ and $\LJIDHA$
are defined to be obtained from $\CLJIDomega$ and $\LJID$
by adding Heyting arithmetic.
Namely, we add constants and function symbols $0, s, +, \times$,
the inductive predicate symbol $\N$, the productions for $\N$,
and Heyting axioms:
\[
\infer{\N0}{}
\qquad
\infer{\N s x}{\N x}
\qquad
\prove \N x \imp s x \ne 0, \qquad
\prove \N x \land \N y \imp s x=s y \imp x=y, \\
\prove \N x \imp x + 0 = x, \qquad
\prove \N x \land \N y \imp x + s y = s(x+y), \\
\prove \N x \imp x \times 0 = 0, \qquad
\prove \N x \land \N y \imp x \times s y = x \times y + x.
\]
\end{Def}

\section{Stage Numbers for Inductive Definitions}\label{sec:stagenumber}

In this section, we define and discuss stage transformation.

We say a path $\pi$ is in a proof $\Pi$ if $\pi$ is syntactically in $\Pi$
without unfolding of $\Pi$. 
For a finite path that goes up from the lowermost sequent $J_2$ to the uppermost sequent $J_1$ in $\Pi$,
we call $J_{2}$ and $J_{1}$ 
the {\em bottom sequent} and the {\em top sequent} of the path respectively.
We write $\pi_1\pi_2$ for the path of the composition of
paths $\pi_1$ and $\pi_2$,
which goes from the bottom sequent of $\pi_1$
to the top sequent of $\pi_2$.

We introduce a stage number to each inductive atomic formula 
so that
the argument of the formula comes into the inductive predicate at the stage 
of the stage number.
This stage number will decrease by a progressing trace.
A proof in $\LJIDHA$ will be constructed by 
using the induction on stage numbers.

First we give stage transformation of an inductive atomic formula.
We assume a fresh inductive predicate symbol $P'$ for
each inductive predicate symbol $P$
and we call it a {\em stage-number inductive predicate symbol}.
$P'(\Vec t,v)$ means that
the element $\Vec t$ comes into $P$ at the stage $v$.
We transform $P(\Vec t)$ into $\exists vP'(\Vec t,v)$.
We call a variable $v$ a {\em stage number}
of $\Vec t$ when $P'(\Vec t,v)$.
$P(\Vec t)$ and $\exists vP'(\Vec t,v)$ will become equivalent 
by inference rules introduced by the transformation of production rules.
We call $P'(\Vec t,v)$ a {\em stage-number inductive atomic formula}.

Secondly 
we give stage transformation of a production rule.
We transform the production of $P$ (the first rule) in Figure \ref{fig:2}
into the production of $P'$ (the second rule) in Figure \ref{fig:2},
where $v,v_1,\ldots,v_m$ are fresh variables.
\begin{figure*}[t]
\[
\infer{P\Vec t}{Q_1\Vec u_1 & \ldots & Q_n\Vec u_n & P_1\Vec t_1 & \ldots & P_m\Vec t_m}
\qquad
\infer{P'\Vec tv}{Q_1\Vec u_1 & \ldots & Q_n\Vec u_n & v > v_1 & P'_1\Vec t_1v_1 & \ldots & v > v_m & P'_m\Vec t_mv_m & \N v}
\]
\caption{Production Rules}\label{fig:2}
\end{figure*}

Next we give stage transformation of a sequent.
For given fresh variables $\Vec v$,
we transform a sequent $J$ into $J^\circ_{\Vec v}$ defined as follows.
We define $\Gamma^\bullet$ 
as the sequent obtained from $\Gamma$
by replacing $P(\Vec t)$ by $\exists vP'(\Vec t,v)$.
%
For fresh variables $\Vec v$, 
we define $(\Gamma)^\circ_{\Vec v}$ 
as the sequent obtained from $\Gamma^\bullet$
by replacing the $i$-th element 
of the form $\exists vP'(\Vec t,v)$ in the sequent $\Gamma^\bullet$
by $P'(\Vec t,v_i)$.
%
We define $(\Gamma \prove \Delta)^\bullet$ by $\Gamma^\bullet \prove \Delta^\bullet$,
and define $(\Gamma \prove \Delta)^\circ_{\Vec v}$ by $(\Gamma)^\circ_{\Vec v} \prove \Delta^\bullet$.

We write $(a_i)_{i \in I}$ for the sequence of elements $a_i$ 
where $i$ varies in $I$.
We extend a proof to that with open assumptions.
We write $\Gamma \prove_\CLJIDHA \Delta$ with assumptions $(J_i)_{i \in I}$
when there is a proof with assumptions $(J_i)_{i \in I}$ and
the conclusion $\Gamma \prove \Delta$ in $\CLJIDHA$.

By the definition of stage transformation of production rules,
$P'\Vec tv$ implies $\N v$.
\begin{Lemma}
(1) $P'\Vec tv \prove_\CLJIDHA \N v$.

(2) $P'\Vec tv \prove_\LJIDHA \N v$.
\end{Lemma}
{\em Proof.}
(1) and (2) are proved by (Case $P'$) and $(\Ind\ P')$ respectively.
$\Box$

\begin{Def}\rm
In a path $\pi$ in a proof,
we define $\Ineq(\pi)$ as the set of the forms
$v > v'$ and $v = v'$ for any stage numbers $v,v'$
eliminated by every case distinction in $\pi$.
\end{Def}
For example, when the path $\pi$ has the case distinction
\[
\Gamma^\circ, u=t(y), \hat v=v,Q_1u_1(y), \ldots, Q_nu_n(y), 
\\ \qquad
v > v_1, P'_1t_1(y)v_1, 
\ldots, v > v_m, P'_mt_m(y)v_m, 
\\ \qquad
\N v \prove \Delta^\bullet
\]
we take $\hat v=v,v>v_1,\ldots,v>v_m$ for $\Ineq(\pi)$.

The proof of the next proposition gives 
{\em stage transformation}
of a proof into
a proof of the stage transformation of the conclusion of the original proof.
We write $\Pi^\circ$ for the stage transformation of $\Pi$.

\begin{Prop}[Stage Transformation]\label{prop:birthtrans1}
For any fresh variables $\Vec v$, 
$\Gamma \prove_\CLJIDHA \Delta$ with assumptions $(\Gamma_i \prove \Delta_i)_{i \in I}$
without any buds
implies 
$(\Gamma)^\circ_{\Vec v} \prove_\CLJIDHA \Delta^\bullet$
with assumptions $(\Ineq(\pi_i),(\Gamma_i)^\circ_{\Vec v_i} \prove \Delta_i^\bullet)_{i \in I}$
without any buds
for some fresh variables $(\Vec v_i)_{i \in I}$,
where $\pi_i$ is the path from the conclusion to the assumption 
$(\Gamma_i)^\circ_{\Vec v_i} \prove \Delta_i^\bullet$.
\end{Prop}

{\em Proof.}
By induction on the proof.
We will transform
the proof $\Pi$ of $\Gamma \prove_\CLJIDHA \Delta$ with assumptions $(J_i)_{i \in I}$ into
the proof of $\Gamma^\circ_{\Vec v} \prove_\CLJIDHA \Delta^\bullet$ with 
assumptions $(\Ineq(\pi_i),(J_i)^\circ_{\Vec v_i})_{i \in I}$ by
transforming each rule as follows.

Case 1. The rule is not ($P_i$ R), (Case), (Cut), 
(Axiom) with a common inductive atomic formula,
or logical rules with
some of the main formula and the auxiliary formulas being
an inductive atomic formula in the antecedent.

Since the main formula and the auxiliary formulas
are not inductive atomic formulas in the antecedent, 
we can just replace each sequent $J$ by $(J)^\circ_{\Vec v}$.
If the rule is an assumption, since $I=\{1\}$,
we take $\Vec v_1$ to be $\Vec v$.

For example,
\[
\infer[(\lor R)]{\Gamma \prove \Delta,P_1t_1 \lor P_2t_2}{
\Gamma \prove \Delta,P_1t_1,P_2t_2
}
\]
is transformed into
\[
\infer[(\lor R)]{(\Gamma)^\circ_{\Vec v} \prove \Delta^\bullet,\exists v_1P'_1t_1v_1 \lor \exists v_2P'_2t_2v_2}{
(\Gamma)^\circ_{\Vec v} \prove \Delta^\bullet,\exists v_1P'_1t_1v_1, \exists v_2P'_2t_2v_2
}
\]

Case 2. The rule is (Axiom) with a common inductive atomic formula.

We transform
\[
\infer[(Axiom)]{\Gamma,P(\Vec t) \prove \Delta,P(\Vec t)}{}
\]
into 
\[
\infer[(\exists R)]{\Gamma^\circ_{\Vec v},P'(\Vec t,v) \prove \Delta^\bullet,\exists vP'(\Vec t,v)}{
	\infer[(Axiom)]{\Gamma^\circ_{\Vec v},P'(\Vec t,v) \prove \Delta^\bullet,P'(\Vec t,v)}{}
}
\]

Case 3. The rule is (Cut),
or logical rules with
some of the main formula and the auxiliary formulas being
an inductive atomic formula in the antecedent.

We replace each sequent $J$ by $(J)^\circ_{\Vec v}$ and
use
\[
\infer[(\exists L)]{\Gamma^\circ_{\Vec v},\exists vP'(\Vec t,v) \prove \Delta^\bullet}{\Gamma^\circ_{\Vec v},P'(\Vec t,v) \prove \Delta^\bullet}
\]

Note that by IH we can assume each $P'(\Vec t,v)$ has a fresh $v$ in $\Gamma^\circ_{\Vec v}$,
so we can use $(\exists L)$.
For example, we transform
\[
\infer[(\neg R)]{\Gamma \prove \neg Pt,\Delta}{\Gamma,Pt \prove \Delta}
\]
into
\[
\infer[(\neg R)]{\Gamma^\circ_{\Vec v} \prove \neg \exists vP'tv,\Delta^\bullet}{
\infer[(\exists L)]{\Gamma^\circ_{\Vec v},\exists vP'tv \prove \Delta^\bullet}{
\Gamma^\circ_{\Vec v},P'tv \prove \Delta^\bullet
}}
\]

Case 4. The rule is $(P\ R)$. 
Assume the production rule of $P$ 
and its stage transformation in Figure \ref{fig:2}.
Let $\Sigma_0$ be the set of all assumptions of this production rule of $P'$
and $t_0$ be $v_1+\ldots+v_n+1$.
By $(P'\ R)$ for this production rule,
\[
\Sigma_0 \prove P'\Vec tv.
\]
Since $P_i\Vec t_iv_i \prove \N v_i$, $(\Sigma_1)^\circ_{\Vec v} \prove \N v_i$
where $\Sigma_1$ is $Q_1\Vec u_1, \ldots, Q_m\Vec u_m,
P_1\Vec t_1, \ldots, P_n\Vec t_n$ and $\Vec v$ is $v_1\ldots v_n$.
By (Subst) with $v:=t_0$,
and
(Cut) with $\prove t_0 > v_i$ and $(\Sigma_1)^\circ_{\Vec v} \prove \N t_0$,
\[
(\Sigma_1)^\circ_{\Vec v} \prove P'\Vec tt_0.
\]
By $(\exists R)$, then by $(\exists L)$,
$
(\Sigma_1)^\bullet \prove \exists vP'\Vec tv.
$

By IH,
\[
\Gamma^\circ_{\Vec v} \prove Q_j\Vec u_j, \Delta^\bullet, \\
\Gamma^\circ_{\Vec v} \prove \exists v_iP'_i\Vec t_iv_i, \Delta^\bullet.
\]
By (Cut),
$
\Gamma^\circ_{\Vec v} \prove \exists vP'\Vec tv, \Delta^\bullet.
$

Case 5. The rule is $(\Case\ P)$.
Assume the production rule of $P$ 
and its stage transformation in Figure \ref{fig:2}.
Let $\Vec v$ be $\Vec v'\hat v$.
Let $v_1,\ldots,v_m$ be fresh variables.
Let the rule be
\[
\infer[(\Case\ P)]{\Gamma,P\Vec u \prove \Delta}{
\hbox{case distinctions}
}
\]
with the case distinctions
\[
\Gamma, \Vec u=\Vec t(\Vec y),Q_1\Vec u_1(\Vec y),\ldots, Q_n\Vec u_n(\Vec y), 
\\ \qquad
P_1\Vec t_1(\Vec y), \ldots, P_m\Vec t_m(\Vec y)
\prove \Delta.
\]

Let
\[
\Sigma \equiv (\Vec u=\Vec t(\Vec y), Q_1\Vec u_1(\Vec y), \ldots, Q_n\Vec u_n(\Vec y),
\\ \qquad
P'_1\Vec t_1(\Vec y)v_1, \ldots, P'_m\Vec t_m(\Vec y)v_m), \\
\Sigma_1 \equiv (\hat v=v,v > v_1,\ldots,v > v_m).
\]

By IH with $\Vec v'v_1\ldots v_m$ for the case distinction
we obtain a proof of
\[
\Gamma^\circ_{\Vec v'}, \Sigma \prove \Delta^\bullet
\]
with assumptions $(\Ineq(\pi'_i),(J_i)^\circ_{\Vec v_i})_{i \in I}$
for some $(\Vec v_i)_{i \in I}$
where 
$\pi'_i$ is the path from the transformation of the case distinction
to $(J_i)^\circ_{\Vec v_i}$.
Add $\Sigma_1$
to the antecedent of every sequent in the path $\pi'_i$,
by renaming fresh variables to keep freshness of all fresh variables
in the path.
Then we have a proof of
\[
\Sigma_1,
\Gamma^\circ_{\Vec v'}, \Sigma \prove \Delta^\bullet
\]
with assumptions $(\Sigma_1,\Ineq(\pi_i'),(J_i)^\circ_{\Vec v_i})_{i \in I}$.
We have
\[
\infer[(\Case\ P')]{\Gamma^\circ_{\Vec v}, P'\Vec u\hat v \prove \Delta^\bullet}{
\hbox{case distinctions}
}
\]
with the case distinctions
\[
\Gamma^\circ_{\Vec v}, \Sigma_1,\Sigma, \N v \prove \Delta^\bullet.
\]
By (Wk) with $\N v$ we obtain the case distinction.
By (Case $P'$),
\[
\Gamma^\circ_{\Vec v'}, P'\Vec u\hat v \prove \Delta^\bullet.
\]
Let $\pi_i$ be the path from $\Gamma^\circ, P'\Vec u\hat v \prove \Delta^\bullet$
to $(J_i)^\circ_{\Vec v_i}$.
Then
\[
\Ineq(\pi_i) = (\Sigma_1,\Ineq(\pi'_i)).
\]
Hence we have a proof of
\[
\Gamma^\circ_{\Vec v'}, P'\Vec u\hat v \prove \Delta^\bullet
\]
with assumptions $(\Ineq(\pi_i),(J_i)^\circ_{\Vec v_i})_{i \in I}$.
$\Box$

\begin{Prop}\label{prop:birthtrans2}
For any fresh variables $\Vec v$,
$\Gamma^\circ_{\Vec v} \prove_\LJIDHA \Delta^\bullet$ 
implies $\Gamma \prove_\LJIDHA \Delta$.
\end{Prop}

{\em Proof.}
First we can show $\Gamma^\bullet \prove_\LJIDHA \Delta^\bullet$
by 
\[
\infer[(\exists L)]{
\exists v P'\Vec tv,\Gamma^\bullet \prove \Delta^\bullet}{
P'\Vec tv,\Gamma^\bullet \prove \Delta^\bullet
}
\]

Secondly
by
\[
\prove_\LJIDHA P\Vec t \lequiv \exists vP'(\Vec t,v)
\]
we have 
\[
\Gamma \prove_\LJIDHA \Delta.
\]
$\Box$

\section{Proof Transformation}\label{sec:prooftrans}

This section defines main proof transformation from $\CLJIDHA$ to $\LJIDHA$.
First we will define path relations, and then we will define main
proof transformation by using path relations.

For notational convenience,
we assume a cyclic proof $\Pi$ in this section.
Let the buds in $\Pi$ be $J_{1i} \ (i \in I)$ and
the companions be $J_{2j} \ (j \in K)$.
Assume $f:N \to N$ such that the companion of a bud $J_{1i}$ is $J_{2,f(i)}$.

\subsection{Path Relation}

In this section, we will introduce path relations and discuss them.

For a path $\pi$,
we will define the path relation $>_\pi$ 
so that it abstracts $\Ineq(\pi)$ and the set of path relations of all paths
is finite.

We assume a subproof $\Pi_j$ of $\Pi$ 
such that it does not have buds,
its conclusion is $J_{2j}$ and its assumptions are
$J_{1i} \ (i \in I_j)$. 

For $J$ in $\Pi^\circ_j$, we define $\Tilde J$ as $\< v_1,\ldots,v_k \>$ where
$J$ is $\Gamma^\circ_{v_1\ldots v_k} \prove \Delta^\bullet$.
Note that 
$\Gamma^\circ$ contains $P'_1\Vec t_1v_1,\ldots,P'_k\Vec t_kv_k$
by definition of $(\ )^\circ$.

For a path $\pi$ from the conclusion to an assumption in $\Pi_j^\circ$,
we write $\check\pi$ for the corresponding path in $\Pi$.
We extend this notation to a finite composition of $\pi$'s.
By the correspondence $(\check{\ })$,
a path, a trace, a progressing trace,
and a stage-number inductive atomic formula in $\Pi^\circ_j$
correspond to those and an inductive atomic formula in $\Pi$.

\begin{Def}\rm
For a finite composition $\pi$ of paths in $\{ \Pi_j^\circ \ |\ j \in K\}$ such that
$\check\pi$ is a path in the infinite unfolding of $\Pi$,
we define the {\em path relation} $\tilde>_{\pi}$ by
\[
x \tilde>_{\pi} y \equiv
|x|=|\Tilde J_2| \land |y|=|\Tilde J_1| \land
\\ \qquad
\Land_{F(q_1,q_2)} (x)_{q_2} > (y)_{q_1}
\land
\Land_{G(q_1,q_2)} (x)_{q_2} = (y)_{q_1}
\]
where 
$J_1$ and $J_2$ are the top and bottom sequents of $\pi$ respectively,
$\check J_1$ and $\check J_2$ are those of the path $\check\pi$,
$F(q_1,q_2)$ is that
there is some progressing trace from the $q_2$-th atomic formula in $\check J_2$
to the $q_1$-th atomic formula in $\check J_1$,
$G(q_1,q_2)$ is that
there is some non-progressing trace from the $q_2$-th atomic formula in $\check J_2$ to the $q_1$-th atomic formula in $\check J_1$.
\end{Def}

\begin{Lemma}\label{lemma:decrease}
For a proof $\Pi$ without any buds,
if $\pi$ is a path from $(J_2)^\circ_{\Vec x}$ to $(J_1)^\circ_{\Vec y}$ in $\Pi^\circ$,
\[
\Ineq(\pi) \prove_\HA \<\Vec x\>  \tilde>_{ \pi} \<\Vec y\>.
\]
\end{Lemma}

{\em Proof.}
By induction on $|\pi|$. 

Case 1. $|\pi|=0$.
$J_1 = J_2$. $x \tilde>_{ \pi} y$ is
$|x|=m \land |y|=m \land (x)_0=(y)_0 \land \ldots \land (x)_{m-1}=(y)_{m-1}$
where $m$ is the number of stage-number inductive atomic formulas in $(J_1)^\circ_{\Vec y}$.
Hence $\prove_\HA x \tilde>_{ \pi} y$.

Case 2. $|\pi|>0$.
Consider cases according to the last rule of $\pi$.
Let $\pi=\pi_1\pi_2$ such that $|\pi_1|=1$.
Let the top sequent of $\pi_1$ be $(J_3)^\circ_{\Vec z}$.
Let $x$,$y$,$z$ be $\<\Vec x\>$, $\<\Vec y\>$, $\<\Vec z\>$ respectively.

By IH, $\Ineq(\pi_2) \prove_\HA z \tilde>_{ \pi_2} y$.

We will show $\Ineq(\pi_1) \prove_\HA x \tilde>_{ \pi_1} z$.
Since the rule that changes the stage number is only (Case),
we will show only the case (Case).
Assume the production rule of $P$ 
and its stage transformation in Figure \ref{fig:2}.
Let the path for the rule (Case $P'$) be
\[
\infer{\Gamma^\circ_{\Vec v}, P'\Vec u\hat v \prove \Delta^\bullet}{
\Gamma^\circ_{\Vec v}, \Sigma_1,\Sigma, \N v \prove \Delta^\bullet
}
\]
where
\[
\Sigma \equiv (\Vec u=\Vec t(\Vec y), Q_1\Vec u_1(\Vec y), \ldots, Q_n\Vec u_n(\Vec y),
\\ \qquad
P'_1\Vec t_1(\Vec y)v_1, \ldots, P'_m\Vec t_m(\Vec y)v_m), \\
\Sigma_1 \equiv (\hat v=v,v > v_1,\ldots,v > v_m).
\]
Then $x=\<\Vec v,\hat v\>$ and $z=\<\Vec v,v_1,\ldots,v_m\>$, and
\[
x \tilde>_{ \pi_1} z \lequiv v_1 < \hat v \land \ldots \land v_m < \hat v.
\]
Since $\Ineq(\pi_1) \equiv \Sigma_1$,
$
\Ineq(\pi_1) \prove_\HA x \tilde>_{ \pi_1} z.
$
Since $\Ineq(\pi) = (\Ineq(\pi_1), \Ineq(\pi_2))$,
$
\Ineq(\pi) \prove_\HA 
x \tilde>_{ \pi_1} z \tilde>_{ \pi_2} y.
$
Since 
$(\tilde>_{ \pi_1}) \circ (\tilde>_{ \pi_2}) \subseteq (\tilde>_{\pi_1\pi_2})$,
$
\Ineq(\pi) \prove_\HA x \tilde>_{ \pi_1\pi_2} y.
$
$\Box$

We define $B_0$ as the set of paths from conclusions to assumptions in 
$\Pi^\circ_j \ (j \in K)$.
We define $B$ as the set of finite compositions of elements in $B_0$ such that
if $\pi \in B$ then $\check\pi$ is a path in the infinite unfolding of $\Pi$.

\begin{Def}\rm\label{def:>}
For $\pi \in B$,
define $x >_{\pi} y$ by
\[
x >_{\pi} y \equiv
(x)_0=j \land (y)_0=f(i) \land  (x)_1 \tilde>_{\pi} (y)_1.
\]
where 
$J_{1i}$ is the top sequent of $\check\pi$, and
$J_{2j}$ is the bottom sequent of $\check\pi$.
\end{Def}

Note that $>_{\pi}$ is a relation on $N \times N^{\le p}$ where
$p$ is the maximum number of inductive atomic formulas in the antecedents of $\Pi$.
The first element is a companion number.

By this note, the next lemma immediately follows.
\begin{Lemma}\label{lemma:finite}
$\{ >_{\pi} \ |\ \pi \in B \}$ is finite.
\end{Lemma}

We define $>_\Pi$ as $\bigcup \{ >_{\pi} \ |\ \pi \in B \}$.
Note that $>_\Pi$ is transitive,
since 
the top sequent of $\pi_1$
is the bottom sequent of $\pi_2$ by the first element, and
$((>_{\pi_1}) \circ (>_{\pi_2})) \subseteq (>_{\pi_1\pi_2})$.

\subsection{Proof Transformation}

This section gives main proof transformation.

First we prepare a piece of main transformation by
proving that we can replace (Case) rules of $\CLJIDHA$ 
by (Ind) rules of $\LJIDHA$ if there are no bud-companion relations.

\begin{Lemma}\label{lemma:case-ind}
If there is a proof
with some assumptions
and without any buds in $\CLJIDHA$,
then
there is a proof of the same conclusions with the same assumptions
in $\LJIDHA$.
\end{Lemma}

{\em Proof.}
It is sufficient to replace the rule (Case) by the rule $(\Ind)$.
This is straightforward and
has been proved by Lemma 4.1.4 in \cite{Brotherston06}.
We give only a proof idea here.

Assume the production rule of $P$ in Figure \ref{fig:2}.
We can replace
{
\[
\infer[(Case)]{\Gamma,P\Vec u \prove \Delta}{
	\infer*[\Pi_i]{\Gamma,\Vec u=\Vec t(\Vec y),\Sigma \prove \Delta}{}
	&
	(\forall i \in I)
}
\]
}
by
{
\[
\infer[(\imp L)(\forall L)(Cut)]{\Gamma,P\Vec u \prove \Delta}{
\infer[(\Ind\ P)]{P\Vec u \prove F\Vec u}{
\infer*[\Pi'_i]{
\Sigma' \prove F\Vec t(\Vec y),F\Vec u}{}
&
\infer{F\Vec u \prove F\Vec u}{}
}}
\]
}
where
\[
F'\Vec z \equiv \forall \Vec x(\Vec z=\Vec u \imp \Gamma \imp \Delta), \\
F\Vec z \equiv P\Vec z \land F'\Vec z.
\]
$\Box$

For $j \in K$, define
\[
G_jx \equiv \forall \Vec z(x=\<\Vec v\> \imp (\Gamma_{2j})^\circ_{\Vec v} \imp \Delta_{2j}^\bullet)
\]
where $\Vec z$ is $\FV((\Gamma_{2j})^\circ_{\Vec v},\Delta_{2j}^\bullet)$.

Define
\[
Gx_0x \equiv \Land_{j \in K} (x_0=j \imp G_jx), \\
Hx_0x \equiv \forall y_0y.
\<y_0,y\> <_\Pi \<x_0,x\> \imp Gy_0y.
\]

\begin{Lemma}\label{lemma:prooftrans}
Assume $\Ind(U,>_\Pi,G)$.
For every bud $J$ of a proof in $\CLJIDHA$ and fresh variables $\Vec v$,
$(J)^\circ_{\Vec v}$ is provable in $\LJIDHA$.
\end{Lemma}

{\em Proof.}
We will show that 
for every companion $J$ in $\Pi$ and fresh variables $\Vec v$,
$(J)^\circ_{\Vec v}$ is provable in $\LJIDHA$.
Fix a companion in $\Pi$ and fresh variables $\Vec v$.
Let the companion be $J_{2j}$.
We will construct a proof of $(J_{2j})^\circ_{\Vec v}$ in $\LJIDHA$.
We have a subproof $\Pi_j$ of $\Pi$ 
such that it does not have buds,
its conclusion is $J_{2j}$ and its assumptions are
$J_{1i} \ (i \in I_j)$. 
By Proposition \ref{prop:birthtrans1},
\[
\infer*[\Pi_j^\circ]{(\Gamma_{2j})^\circ_{\Vec v} \prove \Delta_{2j}^\bullet}{
(\Ineq(\pi_{ji}), (\Gamma_{1i})^\circ_{\Vec v_i} \prove \Delta_{1i}^\bullet)
_{i \in I_j}
}
\]
for some fresh variables $(\Vec v_i)_{i \in I_j}$,
where
$\pi_{ji}$ is the path from $J_{2j}$ to $J_{1i}$ in $\Pi_j$.
Next we transform $\Pi_j^\circ$ 
into the ($\LJIDHA$)-proof $\Pi'_j$ 
with the same conclusion and the same assumptions
by Lemma \ref{lemma:case-ind}.

Next we add
\[
x_0=j,x=\<\Vec v\>,Hx_0x
\]
to every antecedent in $\Pi'_j$ 
to obtain 
\[
\infer*[\Pi''_j]{x_0=j,x=\<\Vec v\>,Hx_0x,(\Gamma_{2j})^\circ_{\Vec v} \prove \Delta^\bullet_{2j}}{
(x_0=j,x=\<\Vec v\>,Hx_0x,\Ineq(\pi_{ji}),(\Gamma_{1i})^\circ_{\Vec v_i} \prove \Delta^\bullet_{1i})
_{i \in I_j}
}
\]

Let $\Pi_{ij}$ be the proof in Figure \ref{fig:3}.
\begin{figure*}[t]
{\tiny
\[
\infer[(\imp L)(\forall L)(Cut)]{x_0=j,x=\<\Vec v\>,Hx_0x,\Ineq(\pi_{ji}),(\Gamma_{1i})^\circ_{\Vec v_i} \prove \Delta^\bullet_{1i}}{
\infer[(Cut)(Subst)]{x_0=j,x=\<\Vec v\>,Hx_0x,\Ineq(\pi_{ji}) \prove Gf(i)\<\Vec v_i\>}{
	\infer[(Wk)(=L)]{x=\<\Vec v\>,y=\<\Vec v_i\>,\Ineq(\pi_{ji}) \prove x \tilde>_{\pi_{ji}} y}{\hbox{Lemma \ref{lemma:decrease}}}
	&
	\infer[(\imp L)(\forall L)(Cut)]{x_0=j,y_0=f(i),
	\<x_0,x\> >_{\pi_{ji}} \<y_0,y\>, Hx_0x \prove Gy_0y}{}
}}
\]
}
\caption{Subproof $\Pi_{ij}$}\label{fig:3}
\end{figure*}
We have a proof $\Pi'''$ of $\forall x_0x.Gx_0x$ 
with the assumption $\Ind(>_\Pi,G)$
in Figure \ref{fig:1}.
\begin{figure*}[t]
\[
\infer[(\imp R)(\forall R)(Cut)]{\prove \forall x_0x.Gx_0x}{
	\prove \Ind(>_\Pi,G)
	&
	\infer[(\land R)]{Hx_0x \prove Gx_0x}{
		(\forall j \in K)
		&
		\infer[(\imp R)(\forall R)(\imp R)(\imp R)]{Hx_0x \prove x_0=j \imp G_jx}{
		\infer*[\Pi''_j]{x_0=j,x=\<\Vec v\>,Hx_0x,(\Gamma_{2j})^\circ_{\Vec v} \prove \Delta^\bullet_{2j}}{
		\left(
		\infer*[\Pi_{ij}]{x_0=j,x=\<\Vec v\>,Hx_0x,\Ineq(\pi_{ji}),(\Gamma_{1i})^\circ_{\Vec v_i} \prove \Delta^\bullet_{1i}}{}
		\right)
		_{i \in I_j}
}}}}
\]
\caption{Proof $\Pi'''$}\label{fig:1}
\end{figure*}

Then
\[
\infer[(\imp L)(\forall L)(Cut)]{(J_{2j})^\circ_{\Vec v}}{
\infer[(\imp L)(\land L)(Cut)]{\prove G_j\<\Vec v\>}{
\infer[(\forall L)(Cut)]{\prove Gj\<\Vec v\>}{
\infer*[\Pi''']{\prove \forall x_0x.Gx_0x}{
\Ind(U,>_\Pi,G)
}}}}
\]

We have shown that 
for every companion $J$ in $\Pi$ and fresh variables $\Vec v$,
$(J)^\circ_{\Vec v}$ is provable in $\LJIDHA$.

Fix a bud be $J_{1k}$ in $\Pi$ and fresh variables $\Vec v$.
Let $J_{2j}$ be the companion of the bud.
Since $(J_{2j})^\circ_{\Vec v}$ is provable, 
$(J_{1k})^\circ_{\Vec v}$ is provable in $\LJIDHA$.
$\Box$

\begin{Prop}[{\tiny Proof Transformation from $\LJIDHA$ to $\CLJIDHA$ Assuming Induction Principle}]
Assume $\Ind(U,>_\Pi,G)$.
$\Gamma \prove_\CLJIDHA \Delta$ with inductive predicates $P_1,\ldots,P_n$
implies $\Gamma \prove_\LJIDHA \Delta$ with $P_1,\ldots,P_n$ and their stage-number inductive predicates $P_1',\ldots,P_n'$.
\end{Prop}

{\em Proof.}
Let $\Pi_1$ be the cyclic proof of $\Gamma \prove_\CLJIDHA \Delta$.
Let $(J_{1i})_{i \in I}$ be all the buds in $\Pi_1$.
Define $\Pi_2$ be a proof obtained from $\Pi_1$ by removing 
all bud-companion relations.
Then
$\Pi_2$ is a proof of $\Gamma \prove \Delta$
with assumptions $(J_{1i})_{i \in I}$ and without buds in $\CLJIDHA$.
Choose fresh variables $\Vec v$.
By Proposition \ref{prop:birthtrans1}, there is $(\Vec v_i)_{i \in I}$ such that
$\Gamma^\circ_{\Vec v} \prove \Delta^\bullet$
with assumptions $(\Ineq(\pi_i),(J_{1i})^\circ_{\Vec v_i})_{i \in I}$ and without buds in $\CLJIDHA$.
By Lemma \ref{lemma:case-ind}, 
$\Gamma^\circ_{\Vec v} \prove \Delta^\bullet$
with assumptions $(\Ineq(\pi_i),(J_{1i})^\circ_{\Vec v_i})_{i \in I}$ and without buds is provable in $\LJIDHA$.
By (Wk),
there is a proof $\Pi_3$ of $\Gamma^\circ_{\Vec v} \prove \Delta^\bullet$
with assumptions $((J_{1i})^\circ_{\Vec v_i})_{i \in I}$ in $\LJIDHA$.
By Lemma \ref{lemma:prooftrans},
$(J_{1i})^\circ_{\Vec v_i}$ is provable in $\LJIDHA$.
Combining them with $\Pi_3$,
$\Gamma^\circ_{\Vec v} \prove \Delta^\bullet$ is provable in $\LJIDHA$.
By Proposition \ref{prop:birthtrans2},
$\Gamma \prove \Delta$ is provable in $\LJIDHA$.

Moreover the proof of the correctness of the proof transformation is
intuitionistic.
$\Box$
